\SourceLeaf{384}%
was much less than it should have been on the assumption of a regular figure of the Earth. The new survey has shown that \textsc{Beccaria} did indeed make an error of \(13''41\) in determining this curvature astronomically (a very pardonable error given the imperfection of his instruments). The sign of this error, however, is opposite to the one suspected, so that the anomaly is increased by that very amount. The amplitude calculated from the geodetic measurements using the elements given above is (reduced to \textsc{Beccaria's} endpoints) \(1^0\,8'\,18''91\), the one derived from \textsc{Beccaria's} astronomical observations is \(1^0\,7'\,44''30\), and the new determination is \(1^0\,7'\,31''07\). The new survey was carried out with such good equipment and with such exceptional care that one is forced to regard this large difference of \(47''84\) almost entirely as an irregularity of the Earth's figure, the most remarkable fact of this kind to have appeared so far in the annals of higher geodesy. The attraction of the Alpine chains that bound this survey to the north and south is most probably a principal cause of this phenomenon, but it is just as probable that the uneven density of the deeper layers of the Earth, perhaps down to a great depth, plays no smaller part. At any rate, similar differences of very considerable magnitude found between points lying entirely on level ground (e.\,g. an anomaly of \(21''9\) between Milan and Parma) can scarcely be explained otherwise. We add that as carefully conducted surveys of arcs become more numerous, the conviction increasingly takes hold that deviations of this kind may be regarded as exceptions only in respect of their magnitude and that they occur everywhere on a larger or smaller scale. The authors have given an interesting comparative survey of the altitudes of the pole determined by astronomical observations and those calculated by geodetic measurements at 34 points scattered over half of Europe and linked by triangles. It must of course be regarded as an imperfect attempt, since it rests for the most part only on unauthenticated, fragmentary notes on the results of the geodetic measurements. Unfortunately, most of these measurements in France, Italy, Austria and Bavaria have still not been published.

This section also contains a new measurement of a short baseline near Turin, which sheds more light on certain circumstances concerning the triangulation carried out there by Mr.~%
